% TOP_PROBLEM: {"id": 2307, "title": "Erdős Problem #790", "category_id": 2, "category": {"id": 2, "name": "combinatorics", "display_name": "Combinatorics", "description": "Counting problems, graph theory, discrete structures.", "slug": "combinatorics", "order_index": 2, "created_at": "2026-07-31T15:26:25.671Z"}, "status": "open"}
% FIRST_POSTED: null
% ENTRY_KIND: statement_only
% Imported from: unsolvedmath_additional_500.tex; UnsolvedMath ID 2307
% !TeX program = lualatex
% Compile twice: lualatex 2307.tex
% Self-contained: no external bibliography, images, or \input files.
% Numeric section labels and visible numbers are original UnsolvedMath IDs.
\documentclass[11pt,a4paper]{article}
\usepackage[margin=24mm,headheight=14pt]{geometry}
\usepackage{fontspec}
\setmainfont{Latin Modern Roman}
\setsansfont{Latin Modern Sans}
\setmonofont{Latin Modern Mono}
\usepackage{amsmath,amssymb,mathtools}
\usepackage{unicode-math}
\setmathfont{Latin Modern Math}
\usepackage{microtype}
\usepackage{enumitem,xurl,needspace}
\usepackage[dvipsnames]{xcolor}
\usepackage{hyperref}
\hypersetup{unicode=true,colorlinks=true,linkcolor=MidnightBlue,urlcolor=MidnightBlue,
 pdftitle={UnsolvedMath Problem 2307},pdfauthor={Source-annotated compilation},bookmarksnumbered=true}
\usepackage{bookmark,tocloft,titlesec,fancyhdr}
\setlength{\cftsecnumwidth}{4.2em}
\cftsetpnumwidth{2.5em}
\cftsetrmarg{4em}
\titleformat{\section}{\Large\bfseries\raggedright}{\thesection}{0.7em}{}
\pagestyle{fancy}\fancyhf{}
\fancyhead[L]{\small\sffamily Additional UnsolvedMath statements}
\fancyhead[R]{\small Original catalogue IDs}
\fancyfoot[C]{\thepage}
\setlength{\parindent}{0pt}
\setlength{\parskip}{5pt plus 1pt}
\setlength{\emergencystretch}{3em}
\setcounter{tocdepth}{1}\setcounter{secnumdepth}{1}
\setlist[itemize]{leftmargin=3em,itemsep=4pt,topsep=4pt,parsep=0pt}
\allowdisplaybreaks
\newcommand{\N}{\mathbb{N}}
\newcommand{\Z}{\mathbb{Z}}
\newcommand{\Q}{\mathbb{Q}}
\newcommand{\R}{\mathbb{R}}
\newcommand{\C}{\mathbb{C}}
\newcommand{\F}{\mathbb{F}}
\newcommand{\A}{\mathbb{A}}
\newcommand{\PP}{\mathbb{P}}
\newcommand{\E}{\mathbb{E}}
\newcommand{\eps}{\varepsilon}
\newcommand{\dd}{\,\mathrm d}
\newcommand{\sref}[2]{\hyperlink{source:#1:#2}{[S#2]}}
\newif\ifshowcatalogue
\showcataloguetrue
\newcommand{\cataloguescope}[1]{\ifshowcatalogue
 \paragraph{Statement recorded in the supplied source.}
 #1\fi}
\begin{document}
% UnsolvedMath ID 2307; source catalogue code EP-790

\setcounter{section}{2306}

\section{Subsets Avoiding Every Distinct-Summand Relation}\label{2307}

{\small\sffamily UnsolvedMath ID 2307 · Combinatorics}

\subsection{Definitions and mathematical statement}

\paragraph{Shared notation.}Unless a section says otherwise, $\N=\{1,2,\ldots\}$,
$\Z$ is the ring of integers, $\Q,\R,\C$ are the rational, real, and complex
fields, $[N]=\{1,\ldots,\lfloor N\rfloor\}$, and $\log$ is the natural
logarithm. For real functions of a parameter tending to infinity,
$f\ll g$ and $f=O(g)$ mean $|f|\leq Cg$ eventually for some fixed $C$;
$f\gg g$ reverses this comparison; $f=o(g)$ means $f/g\to0$;
$f\sim g$ means $f/g\to1$; and $f\asymp g$ means both order bounds.
A subscript on $O$ or $\ll$ specifies allowed constant dependence.
The expression $n^{o(1)}$ in an upper bound means growth smaller than
$n^\epsilon$ for each fixed $\epsilon>0$. Local definitions govern any
exception, finite-field convention, or use of the same symbol for another
object. An asymptotic claim is not automatically an effective algorithm.



Here sum-free'' forbids a selected element from equalling a sum of any number of other distinct selected elements. It is stronger than merely forbidding an equation $a+b=c$. The function is the worst-case guaranteed size over all $n$-element integer sets. \sref{2307}{1} \sref{2307}{2}

\paragraph{Statement recorded in the supplied source.}
Let $l(n)$ be maximal such that if $A\subset\mathbb{Z}$ with $\lvert A\rvert=n$ then there exists a sum-free $B\subseteq A$ with $\lvert B\rvert \geq l(n)$ - that is, $B$ is such that there are no solutions to $ a_1=a_2+\cdots+a_r $ with $a_i\in B$ all distinct.
Estimate $l(n)$. In particular, is it true that $l(n)n^{-1/2}\to \infty$? Is it true that $l(n)< n^{1-c}$ for some $c>0$? \sref{2307}{1}

\paragraph{Scope and interpretation.}

The lower bound recorded in the background already implies growth faster than the square-root scale. The sharp growth and the possible fixed-power upper saving are the remaining issues. \sref{2307}{1}

\paragraph{Residual task reported by the archive.}

The embedded review (2026-08-17) records: Determine whether l(n) is n\textasciicircum{}(1-o(1)) or improve the upper bound. \sref{2307}{1}

\subsection{Short English statement}

How large a subset can always be extracted from an integer set when no selected element may equal a sum of other distinct selected elements?

\subsection{Sources}

{\small

\begin{itemize}

\item[\hypertarget{source:2307:1}{S1}] ULAM.AI, \emph{UnsolvedMath}, supplied \texttt{problems.json}, record 2307 (EP-790), statement and background. The embedded literature-review date is 2026-08-17; that is the archive’s review date, not a fresh certification. Dataset landing page: \url{https://huggingface.co/datasets/ulamai/UnsolvedMath}.

\item[\hypertarget{source:2307:2}{S2}] T. F. Bloom, \emph{Erdős Problems}, Problem 790, \url{https://www.erdosproblems.com/790}. Reference imported from the supplied record; the live problem page was not independently checked for this compilation.

\item[\hypertarget{source:2307:3}{S3}] Choi, S. L. G. and Koml\'os, J. and Szemer\'{e}di, E., On sum-free subsequences. Trans. Amer. Math. Soc. (1975), 307--313. Bibliographic information imported from the supplied record.

\item[\hypertarget{source:2307:4}{S4}] Erd\H{o}s, P., Extremal problems in number theory. Proc. Sympos. Pure Math., Vol. VIII (1965), 181-189. Bibliographic information imported from the supplied record.

\item[\hypertarget{source:2307:5}{S5}] Erd\H{o}s, P., Problems and results on combinatorial number theory. A survey of combinatorial theory (Proc. Internat. Sympos., Colorado State Univ., Fort Collins, Colo., 1971) (1973), 117-138. Bibliographic information imported from the supplied record.

\end{itemize}

}

\label{end:2307}
\end{document}
